\chapter{Introduction}
\label{ch:intro}

\section{A Program That Should Not Fail}
\label{sec:motivating}

Consider the fragment of code in Listing~\ref{lst:sb}.
Two threads run
concurrently on two cores of an ordinary multicore processor.
Two shared
variables, \code{x} and \code{y}, are both initialised to zero.
Each thread
writes~1 to one of them and then reads the other.

\begin{lstlisting}[language=C,caption={The store-buffer idiom. Both threads write one shared variable and then read the other.},label={lst:sb}]
/* initially: x == 0, y == 0 */

/* Thread 0 */              /* Thread 1 */
x = 1;                      y = 1;
r0 = y;                     r1 = x;
\end{lstlisting}

Ask the following question: on termination, can both \code{r0} and \code{r1}
hold the value zero?

The argument that they cannot is short and appears airtight.
The two stores are
distinct events; one of them must occur before the other.
Suppose without loss
of generality that \code{x = 1} occurs first.
Thread~1 reads \code{x} only after
it has written \code{y}, and it wrote \code{y} after \code{x = 1} occurred.
Hence Thread~1's read of \code{x} follows the store to \code{x}, and must return
$1$, so \code{r1} $= 1$. By symmetry the same holds if \code{y = 1} occurs first.
The outcome $\texttt{r0} = \texttt{r1} = 0$ therefore requires each read to have
preceded a store that preceded it, which is a contradiction.

This argument is valid.
It is also, on essentially every processor sold today,
empirically false.
Compiled without synchronisation and run in a loop on a
commodity x86 machine, the outcome $\texttt{r0} = \texttt{r1} = 0$ occurs.
On an
ARM machine it occurs more often, and so do several further outcomes that x86
never produces.

The argument fails not because any step of it is wrong, but because of an
unstated premise: that a store, once executed by a core, is immediately visible
to every other core.
That premise is false, and it is false by deliberate design.
It is discarded in exchange for performance, and the mechanism that discards it
--- the \emph{store buffer} --- is present in every high-performance core built in
the last thirty years.

The consequence is not confined to artificial examples.
Listing~\ref{lst:sb} is
the core of Dekker's algorithm, the first published solution to mutual exclusion.
Executed on hardware that permits the outcome above, Dekker's algorithm admits two
threads into the critical section simultaneously.
A correctness proof that has
stood since 1965 does not transfer to the machine.

\section{Coherence Is Not Enough}
\label{sec:coherence-not-enough}

A natural first response is that the machine in question must have a defective
cache.
It does not.
The behaviour survives a perfectly correct implementation of
cache coherence.

Sarangi draws the distinction in \S12.4.2 and \S12.4.3 of the prescribed text
\cite{sarangi2025basic}, and it is worth stating sharply because much of the
confusion surrounding this topic originates in conflating the two notions.

\emph{Coherence} is a guarantee about a single memory location. It requires that
all cores observe the writes to any one address in a single agreed order, and that
a read of that address returns the most recent value in that order.
Protocols such
as MSI, MESI and MOESI, treated in \S12.4.6 of the text and in Chapter~3 of this
paper, exist to provide exactly this.
They are well understood and correctly
implemented in practice.

\emph{Consistency} is a guarantee about the relative order of accesses to
\emph{different} locations. Listing~\ref{lst:sb} involves two addresses. Coherence
constrains the writes to \code{x} among themselves and the writes to \code{y}
among themselves; it says nothing whatever about how Thread~0's store to \code{x}
is ordered with respect to Thread~0's load of \code{y}.
The anomalous outcome
requires no violation of coherence and does not produce one.

\begin{definition}[Memory consistency model]
A \emph{memory consistency model} is the specification, forming part of a
machine's architecture, of which orderings among memory operations issued by a
single core are guaranteed to be preserved as seen by other cores, together with
the instructions provided to the programmer for enforcing orderings that are not
guaranteed by default.
\end{definition}

The model is a contract.
Like the instruction set, it is an interface: hardware
promises no less, software may assume no more.
Unlike the instruction set, it is
frequently absent from a programmer's mental model of the machine entirely, which
is precisely why the failures it permits are so difficult to diagnose.

\section{The Design Tension}
\label{sec:tension}

Figure~\ref{fig:tension} states the situation that this paper investigates.
Two objectives are in conflict.

\begin{figure}[H]
\centering
\begin{tikzpicture}[node distance=0mm]
  % axis
  \draw[flow] (0,0) -- (12.6,0);
  \node[lbl, anchor=west] at (12.7,0) {};

  % ticks
  \foreach \x/\name/\ex in {%
     0.8/{Sequential\\consistency}/{(no shipping ISA)},
     4.2/{TSO}/{x86, SPARC},
     7.2/{PSO}/{SPARC},
    10.8/{Weak / RMO}/{ARM, POWER, RISC-V}}
  {
    \draw (\x,0.12) -- (\x,-0.12);
    \node[align=center, font=\small\bfseries, anchor=south] at (\x,0.22) {\name};
    \node[align=center, font=\scriptsize\itshape, anchor=north] at (\x,-0.22) {\ex};
  }

  % annotations
  \node[align=center, font=\small, anchor=north] at (0.8,-1.3)
        {easy to reason about\\hard to make fast};
  \node[align=center, font=\small, anchor=north] at (10.8,-1.3)
        {fast\\hard to program correctly};

  \draw[decorate,decoration={brace,amplitude=4pt,mirror}]
        (0.2,-2.35) -- (11.4,-2.35);
  \node[lbl, anchor=north] at (5.8,-2.5)
        {increasingly weak contract $\longrightarrow$ more reordering permitted
         $\longrightarrow$ more burden on the programmer};
\end{tikzpicture}
\caption{The design space of hardware memory consistency models. Every
commercially deployed instruction set sits to the right of sequential
consistency.}
\label{fig:tension}
\end{figure}

At the left-hand end of the spectrum the hardware guarantees everything, and the
reasoning of \S\ref{sec:motivating} is sound.
The programmer needs no special
instructions and no special knowledge.
The cost is that the core must, in the
naive implementation, stall on every store until that store is globally visible
--- which is to say, it must abandon the store buffer, and with it a substantial
fraction of its performance.

At the right-hand end the hardware guarantees almost nothing and runs at full
speed. The programmer is handed \emph{fence} instructions and made responsible for
placing them wherever an ordering is genuinely required. Empirically, programmers
place them incorrectly.
The resulting defects are probabilistic, appear only under
particular thread placements and core counts, and typically vanish under a
debugger.
Two decades of experience with the Java Memory Model
\cite{manson2005java}, the C++ memory model \cite{boehm2008foundations} and the
Linux kernel's own memory model attest to the difficulty.

The industry chose the right-hand end.
The justification was performance, and it
was articulated in an era whose assumptions no longer hold.

\section{Why the Question Is Open}
\label{sec:why-open}

Three developments since 1990 bear directly on the original trade-off.

\paragraph{Speculation is now ubiquitous.} A modern out-of-order core executes
instructions ahead of their architectural order and discards the results when a
prediction turns out to have been wrong.
It therefore already possesses a
squash-and-restart mechanism, built for branch misprediction recovery, and it
already receives coherence invalidation messages announcing remote writes.
Together these are sufficient to enforce sequential consistency
\emph{optimistically}: execute the load early, and roll back only in the rare case
where a remote write proves that doing so was observable.
Gniady, Falsafi and
Vijaykumar posed this directly \cite{gniady1999sc}, and later work extended it to
speculative store retirement and to chunk-based execution
\cite{ceze2007bulksc,blundell2009invisifence}.

\paragraph{Compilers matter as much as hardware.} Hardware sequential consistency
is worthless if the compiler reorders shared accesses first.
For a long time it
was assumed that constraining the compiler would be prohibitively expensive.
Marino et al.\ measured it and found the cost modest \cite{marino2011sc}.

\paragraph{Core counts have grown by an order of magnitude.} The measurements that
settled the original debate were taken on machines with between two and sixteen
processors.
Speculation-based SC depends on conflicts being rare; whether that
premise survives at 64 or 128 cores is not established by any of the classical
results.

Against all of this stands one stubborn observation: \emph{no shipping processor
implements sequential consistency}.
If the performance objection has been answered,
something else is keeping the industry where it is --- energy, verification cost,
the irreversibility of an architectural commitment, or simple inertia.

\section{Problem Statement}
\label{sec:problem}

\begin{quote}
Sequential consistency was rejected as a hardware memory model on performance
grounds in the late 1980s, on machines that did not speculate.
This paper asks
whether that verdict still holds.
It (i)~formalises the design space between SC
and the weak models actually deployed within a single ordering framework;
(ii)~establishes by direct measurement on x86 and ARM hardware which relaxations
are observable in practice, as distinct from those merely permitted by the
architecture; and (iii)~quantifies, in a configurable multicore simulator, the
execution-time cost of enforcing each model both with and without speculative
load execution, across core counts beyond those examined in the classical
literature.
\end{quote}

\noindent The problem decomposes into four questions.

\begin{description}
\item[Q1 (Formal).] Can SC, TSO, PSO, weak ordering and release consistency be
  expressed as successively weaker constraints on a single set of relations over
  memory events, and does the resulting hierarchy admit a proof of strict
  containment?
Chapter~4.

\item[Q2 (Empirical).] For a catalogue of litmus tests, which outcomes are
  permitted by each architecture's specification, and which are actually produced
  by hardware?
Where the two differ, what accounts for the difference?
  Chapters~6--7.

\item[Q3 (Quantitative).] Holding the microarchitecture fixed and varying only the
  memory model, what is the execution-time cost of SC relative to TSO and weak
  ordering, with and without speculative load execution, as a function of core
  count?
Chapters~8--9.

\item[Q4 (Interpretive).] If the measured cost of speculative SC is small, what
  accounts for its absence from shipping hardware?
Chapter~11.
\end{description}

\section{Contributions}
\label{sec:contributions}

This paper claims the following.

\begin{enumerate}
\item A unified presentation of five memory consistency models within the
  execution-witness formalism of Sarangi \cite{sarangi2023nextgen}, with the
  containment hierarchy proved rather than asserted.

\item A litmus-test suite, executed on both x86 and ARM hardware, reporting
  observation frequencies rather than the usual binary
  permitted/forbidden classification --- together with a measured cost, in
  cycles, for the fence instructions that suppress each relaxation.

\item A multicore simulator in which the memory consistency model is an
  interchangeable component, permitting a controlled comparison in which the
  model is the only variable.

\item A verification chapter in which the simulator's exhaustively enumerated
  outcomes are checked against the axiomatic definitions of Chapter~4, so that
  the performance results rest on an implementation demonstrated to implement
  the model it claims.

\item A scaling study extending the speculative-SC question to core counts beyond
  those examined in the 1999--2011 literature.
\end{enumerate}

\section{Relation to the Prescribed Text}
\label{sec:relation}

The paper is anchored in Chapter~12 of Sarangi \cite{sarangi2025basic}, and
specifically in \S12.4.
Table~\ref{tab:coverage} maps each chapter of this paper
onto the sections of the text it extends.

\begin{table}[H]
\centering
\small
\caption{Coverage of the prescribed text.}
\label{tab:coverage}
\begin{tabular}{@{}llp{6.1cm}@{}}
\toprule
\textbf{This paper} & \textbf{Text} & \textbf{Relationship} \\
\midrule
Ch.~2 Background   & \S11.2, \S11.3, \S10.11.4, \S12.2 &
  Caches, the memory system, out-of-order pipelines, shared memory versus
  message passing. \\
Ch.~3 Coherence    & \S12.4.2, \S12.4.5, \S12.4.6, \S12.6 &
  Extends the text's treatment of coherent private caches to full MESI and
  MOESI transition tables and to directory protocols. \\
Ch.~4 Models       & \S12.4.3, \S12.4.7 &
  The core section of the text. Extended to a formal treatment of four further
  models and a containment proof. \\
Ch.~5 History      & \S12.7.2 &
  Develops the further-reading pointers into a chronology. \\
Ch.~6 Methodology  & \S12.4.3, \S10.9 &
  Applies the text's performance-metric framework to litmus testing. \\
Ch.~8--9 Simulation & \S10.9, \S12.4.7, \S12.6 &
  Performance equation, consistency implementation, interconnect modelling. \\
Ch.~11 Discussion  & \S12.2.3, App.~A &
  Amdahl's law; the Cortex-A8, Cortex-A15 and Sandy Bridge case studies. \\
\bottomrule
\end{tabular}
\end{table}

Where the treatment exceeds the scope of the prescribed text --- notably the
axiomatic formalism, the relaxed models beyond TSO, and the speculative
implementation techniques --- the primary sources are Chapter~9 of Sarangi's
companion volume \cite{sarangi2023nextgen} and the memory-consistency primer of
Nagarajan et al.\ \cite{nagarajan2020primer}.

\section{Organisation of the Paper}
\label{sec:organisation}

Chapter~2 develops the microarchitectural background: why memory is slow, what
caches and store buffers do about it, and how out-of-order execution compounds
the reordering they introduce.
Chapter~3 treats cache coherence and establishes
precisely what it does and does not guarantee.
Chapter~4 is the formal core:
it defines the ordering relations and expresses each consistency model as a
constraint over them.
Chapter~5 traces the historical argument from Lamport's
1979 paper to the ratification of RVWMO. Chapter~6 sets out the litmus-test
methodology and experimental design for the results reported in later versions
of this paper.

\medskip
\noindent\emph{Scope of version 0.1.} The present version comprises Chapters~1--6.
Chapters~7--12, containing the hardware measurements, the simulator and its
results, the verification study and the discussion, are the subject of versions
v0.2 and v0.3.
